\section*{الفصل \ref{ch:b1:organomagnesium} --- الكواشف العضوية المغنيزيومية}

\begin{solution}{exo:b1:organomagnesium:1}
\ce{CH3CH2Br + Mg -> CH3CH2MgBr}؛ \ce{C6H5Br + Mg -> C6H5MgBr} (في إيثر
جاف). C--Br: البروم (2.96) أكثر كهرسلبيةً من الكربون (2.55)، فالكربون
موجب. C--Mg: المغنيزيوم (1.31) أقل كهرسلبيةً، فالكربون سالب.
\end{solution}

\begin{solution}{exo:b1:organomagnesium:2}
الماء: \ce{CH3MgI + H2O -> CH4 + Mg(OH)I}.
الإيثانول: \ce{CH3MgI + CH3CH2OH -> CH4 + CH3CH2OMgI}.
حمض الإيثانويك: \ce{CH3MgI + CH3COOH -> CH4 + CH3COOMgI}.
أكسيد الدوتيريوم: \ce{CH3MgI + D2O -> CH3D + Mg(OD)I}.
\end{solution}

\begin{solution}{exo:b1:organomagnesium:3}
الميثانال: بروبان-1-أول (أولي). الإيثانال: بوتان-2-أول (ثانوي).
البروبانون: 2-ميثيل بوتان-2-أول (ثالثي). البنزالدهيد:
1-فينيل بروبان-1-أول (ثانوي). ثنائي أكسيد الكربون: حمض البروبانويك.
\end{solution}

\begin{solution}{exo:b1:organomagnesium:4}
مجموعة OH فيه حمضية بما يكفي لهدم الكاشف ما إن يتكوّن:
فالجزيء سيبرتن نفسه (أو جيرانه)، فيعطي
ألكوكسيد المغنيزيوم للبوتان-1-أول، \ce{CH3(CH2)3OMgBr}، لا كاشفًا قابلًا
للاستعمال. ويجب أولًا حماية مجموعة OH
(يبيّن \cref{ch:b1:acetals-protection} الفكرة في حالة الألدهيد).
\end{solution}

\begin{solution}{exo:b1:organomagnesium:5}
يهاجم زوج C--Mg كربون الكربونيل في الإيثانال، ويذهب زوج $\pi$ للرابطة C=O
إلى O: \ce{CH3CH(CH3)O^-} مع \ce{MgBr+}؛ ثم يعطي \ce{H3O+}
بروتونًا للأكسجين: بروبان-2-أول. ويحمل بروبان-2-أول مجموعتي ميثيل
على الكربون الحامل لمجموعة OH: فهو \omterm{def:b1:stereochemistry-in-depth:stereocentre}{لاكيرالي}. (ومع مجموعة R مختلفة، مثل
بروميد إيثيل المغنيزيوم مع الإيثانال، يكون الناتج بوتان-2-أول \omterm{def:b1:stereochemistry-in-depth:stereocentre}{كيراليًا} لكنه
راسيمي: إذ يُهاجَم الكربونيل المستوي على وجهيه بالتساوي.)
\end{solution}

\begin{solution}{exo:b1:organomagnesium:6}
يحمل الكربون الحامل لمجموعة OH مجموعة ميثيل ومجموعتي إيثيل:
بروميد إيثيل المغنيزيوم مع البوتان-2-ون، أو بروميد ميثيل المغنيزيوم مع
البنتان-3-ون.
\end{solution}

\begin{solution}{exo:b1:organomagnesium:7}
حضّر \ce{(CH3)2CHMgBr} من 2-بروموبروبان والمغنيزيوم في إيثر جاف،
واسكبه على \ce{CO2} الصلب، ثم حمّض: \ce{(CH3)2CHCOOH}. $n = 12.3/123.0 = 0.100$ مول؛
$0.100 \times 0.70 \times 88.0 = \qty{6.2}{g}$.
\end{solution}

\begin{solution}{exo:b1:organomagnesium:8}
$0.18/18.0 = \qty{0.010}{mol}$ من الماء تهدم \qty{0.010}{mol} من
الكاشف: \qty{20}{\percent}؛ ويتكوّن البنزين، \ce{C6H5MgBr + H2O -> C6H6 +
Mg(OH)Br}.
\end{solution}

\begin{solution}{exo:b1:organomagnesium:9}
يتفاعل الكاشف، وهو \omterm{def:b1:electronic-effects:nucleophile}{محب للنواة}، مع البروموبنزين الذي ما زال موجودًا:
إجمالًا \ce{C6H5MgBr + C6H5Br -> C6H5-C6H5 + MgBr2}. وإضافة البروميد
ببطء إلى فائض من المغنيزيوم تُبقي تركيزه منخفضًا: فيلقى فلزًا
جديدًا بدل الكاشف المتكوّن من قبل.
\end{solution}

\begin{solution}{exo:b1:organomagnesium:10}
1-فينيل إيثان-1-أول: \ce{CH3MgBr} (من البروموميثان) مع البنزالدهيد،
أو \ce{C6H5MgBr} (من البروموبنزين) مع الإيثانال. 2-فينيل بروبان-2-أول:
\ce{C6H5MgBr} مع البروبانون، أو \ce{CH3MgBr} مع
1-فينيل إيثان-1-ون.
\end{solution}

\begin{solution}{exo:b1:organomagnesium:11}
يُضاف الكيتون إلى الكاشف كي لا يغادر الكاشفُ، المحضَّر في
الدورق الجاف، ذلك الدورقَ أبدًا، وكي يمتص الإيثرُ المغلي الحرارةَ المنطلقة من كل جزء.
والماء وحده كان سيعطي الملح القاعدي الهلامي
\ce{Mg(OH)Br}، الذي يحتجز الناتج؛ أما الحمض المخفف البارد فيذيبه
على هيئة \ce{Mg^2+}، والبرودة تحد من التفاعلات الجانبية
للكحول الثالثي في الوسط الحمضي.
\end{solution}

\begin{solution}{exo:b1:organomagnesium:12}
كل مول من الكاشف يعطي مع الماء مولًا واحدًا من ملح المغنيزيوم القاعدي،
\ce{RMgBr + H2O -> RH + Mg(OH)Br}، الذي يُعاير أيونه \ce{OH-} بالحمض:
فكمية الكاشف تساوي كمية الحمض المستعملة، \qty{0.088}{mol}، أي \omterm{def:b1:extent-q-and-k:fractional-extent}{مردود}
\qty{88}{\percent}.
\end{solution}

\begin{solution}{pb:b1:organomagnesium:1}
\textbf{1.} \ce{C6H5Br + Mg -> C6H5MgBr}.
\textbf{2.} Mg $0.535/24.3 = \qty{0.0220}{mol}$؛ \ce{C6H5Br} $3.14/156.9 =
\qty{0.0200}{mol}$: المغنيزيوم بفائض طفيف.
\textbf{3.} الماء يهدم الكاشف: \ce{C6H5MgBr + H2O -> C6H6 +
Mg(OH)Br}.
\textbf{4.} يمنع دخول بخار ماء الهواء عبر
المكثف.
\textbf{5.} يعطي الإيثر أزواجًا حرة للمغنيزيوم (\omterm{def:b1:lewis-resonance-vsepr:lewis-acid}{قاعدة لويس}):
ومن دونه لا يتكوّن الكاشف ولا يبقى في المحلول.
\textbf{6.} \ce{C6H5MgBr + C6H5Br -> C6H5-C6H5 + MgBr2}؛ ثنائي الفينيل
$0.15/154.0 = \qty{9.7e-4}{mol}$، استهلك \qty{9.7e-4}{mol} من
البروموبنزين ومثلها من الكاشف: فالمفقود من البروموبنزين \qty{1.9e-3}{mol}
إجمالًا.
\textbf{7.} $3.28/182.0 = \qty{0.0180}{mol}$.
\textbf{8.} يهاجم زوج C--Mg كربون الكربونيل، ويذهب زوج $\pi$
إلى الأكسجين: \ce{(C6H5)3C-O^-} \ce{MgBr+}.
\textbf{9.} الكاشف المتاح على الأكثر $0.0200 - 0.0019 = \qty{0.0181}{mol}$؛
والبنزوفينون \qty{0.0180}{mol}: فالبنزوفينون هو المحدِّد، بفارق ضئيل.
\textbf{10.} لا: الكربون الحامل لمجموعة OH يحمل ثلاث مجموعات فينيل متماثلة.
\textbf{11.} يبقى الكاشف في دورقه الجاف، ويمتص الإيثر المغلي حرارة كل
جزء من الكيتون.
\textbf{12.} كان سيهدم جزءًا من الكاشف (بنزين) ويترك
بنزوفينون لم يتفاعل.
\textbf{13.} \ce{(C6H5)3COMgBr + H3O+ -> (C6H5)3COH + Mg^2+ + Br- + H2O}.
\textbf{14.} يذيب الحمض أملاح المغنيزيوم القاعدية التي كان الماء
وحده سيرسّبها؛ والبرودة تحد من تكوّن
\omterm{def:b1:electronic-effects:intermediates}{الكاتيون الكربوني} في السؤال 18.
\textbf{15.} ثلاثي فينيل الميثانول وثنائي الفينيل في طبقة الإيثر؛
وأملاح المغنيزيوم في الطبقة المائية.
\textbf{16.} غسل الصلب (أو سحقه) بإيثر البترول الخفيف البارد:
فيذوب ثنائي الفينيل، ويبقى ثلاثي فينيل الميثانول.
\textbf{17.} درجة انصهار حادة مساوية للقيمة المجدولة (أو بقعة
وحيدة في كروماتوغرافيا الطبقة الرقيقة).
\textbf{18.} \ce{(C6H5)3COH + H+ -> (C6H5)3C+ + H2O}: \omterm{def:b1:electronic-effects:intermediates}{كاتيون كربوني} ثالثي
تنتشر شحنته على ثلاث حلقات بنزين بالرنين، فهو شديد
الثبات، وملوّن بسبب ذلك الترافق الممتد.
\textbf{19.} $3.75/260.0 = \qty{0.0144}{mol}$.
\textbf{20.} $0.0144/0.0180 = \qty{80}{\percent}$.
\textbf{21.} فقد أثناء النقل وإعادة التبلور؛ وكاشف هدمته
آثار الرطوبة؛ وتفاعل غير تام.
\textbf{22.} حمض البنزويك، \ce{C6H5COOH}، بعد التحميض.
\textbf{23.} $0.0181 \times 122.0 = \qty{2.21}{g}$.
\textbf{24.} \ce{C6H5MgBr + CO2 -> C6H5COOMgBr}، ثم
\ce{C6H5COOMgBr + H3O+ -> C6H5COOH + Mg^2+ + Br- + H2O}.
\textbf{25.} \textbf{\qty{80}{\percent}}.
\end{solution}
